\chapter{Boolean Analysis — A Cp Gates}
%%% generated by the blueprint pipeline from TCSlib.BooleanAnalysis.RazborovSmolensky.ACpGates
%%%%%%%%%%%%%%%%%%%%%%%%%%%%%%%%%%%%%%%%%%%%%%%%%%%%%%%%%%%%%%%%%%%%%%%%%%%%%%%
\section{Overview}

This file sets up the gate sets $\mathrm{AC}^0$ and $\mathrm{AC}^0[p]$ over the Boolean
alphabet $\mathrm{Fin}\,2$, and the polynomial machinery used in the Razborov--Smolensky
argument: Razborov's randomized low-degree approximations of unbounded fan-in $\mathrm{OR}$
and $\mathrm{AND}$ over the field $\mathbb{Z}/p$, together with the exact $\mathrm{MOD}_p$
polynomial. The main outputs are total-degree bounds of the form $(p-1)\ell$ per gate and
pointwise counting bounds showing that, for each fixed input, at most a $2^{-\ell}$
fraction of the random seeds produce a wrong value.

%%%%%%%%%%%%%%%%%%%%%%%%%%%%%%%%%%%%%%%%%%%%%%%%%%%%%%%%%%%%%%%%%%%%%%%%%%%%%%%
\section{Declarations}

\begin{definition}[The $\mathrm{AC}^0$ gate set]
\lean{ACP.AC_GateOps}
\label{ACP.AC_GateOps}
\leanok
\uses{ACP.FeedForward.GateOp.id}
The set of plain $\mathrm{AC}^0$ gate operations on the alphabet $\mathrm{Fin}\,2$: the
identity gate, the NOT gate $x \mapsto 1 - x_0$, and, for every arity $n$, the unbounded
fan-in AND gate $x \mapsto \prod_{i} x_i$.
\end{definition}

\begin{lemma}[Counting tuples satisfying a pointwise predicate]
\lean{ACP.tuple_fail_count}
\label{ACP.tuple_fail_count}
\leanok
\difficulty{4}
Let $\iota$ be a finite index type, let $\beta$ be a finite type, and let $P$ be a
decidable predicate on $\beta$. Then the number of functions $f : \iota \to \beta$ such
that $P(f(i))$ holds for every $i \in \iota$ equals $\abs{\{\,b \in \beta :
P(b)\,\}}^{\abs{\iota}}$.
\end{lemma}

\begin{lemma}[Averaging for the probabilistic method]
\lean{ACP.prob_method_averaging}
\label{ACP.prob_method_averaging}
\leanok
\difficulty{5}
Let $\alpha$ be a finite type and $\beta$ a finite nonempty type, let $\mathrm{Bad}
\subseteq \alpha$ be a finite set, let $\mathrm{Fail} \subseteq \alpha \times \beta$ be
a decidable relation, and let $C$ be a natural number. Suppose that for every $a \in
\mathrm{Bad}$ we have $\abs{\{b \in \beta : \mathrm{Fail}(a,b)\}} \cdot C \le
\abs{\beta}$. Then there exists a single $b \in \beta$ such that $\abs{\{a \in
\mathrm{Bad} : \mathrm{Fail}(a,b)\}} \cdot C \le \abs{\mathrm{Bad}}$.
\end{lemma}

\begin{definition}[Booleanization of a field element]
\lean{ACP.bitify}
\label{ACP.bitify}
\leanok
The map $\mathbb{Z}/p \to \mathrm{Fin}\,2$ sending $a$ to $1$ if $a = 1$ and to $0$
otherwise. It is intended to be applied only to values lying in $\{0,1\}$.
\end{definition}

\begin{lemma}[Booleanization is a section on $\{0,1\}$]
\lean{ACP.cast_bitify_eq}
\label{ACP.cast_bitify_eq}
\leanok
\uses{ACP.bitify}
\difficulty{2}
Let $p$ be a prime and let $a \in \mathbb{Z}/p$ satisfy $a \in \{0, 1\}$. Write
$\beta(a) \in \mathrm{Fin}\,2$ for the booleanization of $a$, equal to $1$ when $a = 1$
and to $0$ otherwise. Then interpreting $\beta(a)$ as a natural number and reducing it
back modulo $p$ recovers $a$; that is, the image of $\beta(a)$ under $\mathrm{Fin}\,2
\to \mathbb{N} \to \mathbb{Z}/p$ equals $a$.
\end{lemma}

\begin{lemma}[Fermat indicator of zero mod p]
\lean{ACP.one_sub_pow_card_sub_one}
\label{ACP.one_sub_pow_card_sub_one}
\leanok
\difficulty{3}
Let $p$ be a prime and let $a \in \bbz/p\bbz$. Then
\[
1 - a^{p-1} = \begin{cases} 1 & \text{if } a = 0,\\ 0 & \text{if } a \neq 0.\end{cases}
\]
That is, $1 - a^{p-1}$ is the indicator of $a = 0$ in the field $\bbz/p\bbz$.
\end{lemma}

\begin{lemma}[Fermat indicator computes the boolean value]
\lean{ACP.bit_indicator_eq_bitify}
\label{ACP.bit_indicator_eq_bitify}
\leanok
\uses{ACP.bitify}
\difficulty{2}
Let $p$ be a prime and let $a \in \mathbb{Z}/p$ lie in $\{0, 1\}$. Then
\[
1 - (1 - a)^{p-1} = \mathrm{bitify}(a),
\]
where the right-hand side denotes the element of $\{0,1\}\subseteq \mathbb{Z}/p$
obtained by casting the value $\mathrm{bitify}(a) \in \mathrm{Fin}\,2$ through the
natural numbers into $\mathbb{Z}/p$; that is, the expression $1 - (1 - a)^{p-1}$ equals
$1$ when $a = 1$ and $0$ when $a = 0$.
\end{lemma}

\begin{definition}[Unbounded $\mathrm{MOD}_p$ gate]
\lean{ACP.modGateOp}
\label{ACP.modGateOp}
\leanok
The gate operation of arity \texttt{width} on Boolean inputs that outputs $1$ when the sum
of its inputs, cast into $\mathbb{Z}/p$, is $0$, and outputs $0$ otherwise.
\end{definition}

\begin{definition}[{The $\mathrm{AC}^0[p]$ gate set}]
\lean{ACP.ACp_GateOps}
\label{ACP.ACp_GateOps}
\leanok
\uses{ACP.AC_GateOps, ACP.modGateOp}
The $\mathrm{AC}^0[p]$ gate set: the $\mathrm{AC}^0$ gates (identity, NOT, unbounded AND)
together with the unbounded $\mathrm{MOD}_p$ gate of every arity.
\end{definition}

\begin{definition}[Randomized OR-approximating polynomial]
\lean{ACP.approxOr}
\label{ACP.approxOr}
\leanok
Given polynomials $P_1,\dots,P_{\mathrm{width}}$ over $\mathbb{Z}/p$ and a random seed
consisting of $\ell$ subsets $S_1,\dots,S_\ell$ of the input positions, the approximator is
\[
1 - \prod_{k=1}^{\ell}\Bigl(1 - \bigl(\textstyle\sum_{i \in S_k} P_i\bigr)^{p-1}\Bigr).
\]
\end{definition}

\begin{definition}[Exact $\mathrm{MOD}_p$ polynomial]
\lean{ACP.exactMod}
\label{ACP.exactMod}
\leanok
The polynomial $1 - \bigl(\sum_i P_i\bigr)^{p-1}$ over $\mathbb{Z}/p$, which computes the
$\mathrm{MOD}_p$ predicate exactly on Boolean values.
\end{definition}

\begin{definition}[Value-level OR approximator]
\lean{ACP.approxOr_val}
\label{ACP.approxOr_val}
\leanok
The value-level analogue of the randomized OR approximator: for $v : \mathrm{Fin}\,
\mathrm{width} \to \mathbb{Z}/p$ and subsets $S_1,\dots,S_\ell$,
\[
1 - \prod_{k}\Bigl(1 - \bigl(\textstyle\sum_{i \in S_k} v_i\bigr)^{p-1}\Bigr).
\]
\end{definition}

\begin{definition}[Value-level OR detector]
\lean{ACP.OR_val}
\label{ACP.OR_val}
\leanok
For $v : \mathrm{Fin}\,\mathrm{width} \to \mathbb{Z}/p$, the quantity
$1 - \prod_{k}\bigl(1 - v_k^{\,p-1}\bigr)$, which detects whether some coordinate of $v$ is
nonzero.
\end{definition}

\begin{theorem}[Degree bound for the OR-approximating polynomial]
\lean{ACP.approxOr_totalDegree}
\label{ACP.approxOr_totalDegree}
\leanok
\uses{ACP.approxOr}
\difficulty{4}
Let $p$ be a prime, and work over the ring of polynomials in $\mathrm{vars}$ variables
with coefficients in $\bbz/p$. Given a family of polynomials
$P_1,\dots,P_{\mathrm{width}}$ from this ring, together with a seed consisting of $\ell$
subsets $S_1,\dots,S_\ell$ of the index set $\{1,\dots,\mathrm{width}\}$, the associated
OR-approximating polynomial has total degree at most
\[
(p-1)\,\ell\cdot \sup_i \deg(P_i),
\]
where the supremum is taken over all indices $i$.
\end{theorem}

\begin{theorem}[Total degree of the exact $\mathrm{MOD}_p$ polynomial]
\lean{ACP.exactMod_totalDegree}
\label{ACP.exactMod_totalDegree}
\leanok
\uses{ACP.exactMod}
\difficulty{3}
Let $p$ be a prime and fix natural numbers $\mathrm{vars}$ and $\mathrm{width}$. Let
$P_0,\dots,P_{\mathrm{width}-1}$ be polynomials in $\mathrm{vars}$ variables over the
field $\bbz/p$. Then the total degree of the exact $\mathrm{MOD}_p$ polynomial $1 -
\bigl(\sum_i P_i\bigr)^{p-1}$ is at most
\[
(p-1)\cdot \sup_i \deg(P_i),
\]
the supremum being taken over all $i$ (and read as $0$ when $\mathrm{width}=0$).
\end{theorem}

\begin{lemma}[At most half of subset sums vanish]
\lean{ACP.subset_sum_zero_bound}
\label{ACP.subset_sum_zero_bound}
\leanok
\difficulty{5}
Let $p$ be a prime and let $v = (v_i)_{i \in \{1,\dots,n\}}$ be a nonzero tuple of
elements of the field $\bbz/p\bbz$. Then at most half of the $2^n$ subsets of the index
set $\{1,\dots,n\}$ have vanishing subset sum; that is, twice the number of subsets $s
\subseteq \{1,\dots,n\}$ for which $\sum_{i \in s} v_i = 0$ is at most $2^n$.
\end{lemma}

\begin{lemma}[Evaluation commutes with the OR approximator]
\lean{ACP.approxOr_eval_eq}
\label{ACP.approxOr_eval_eq}
\leanok
\uses{ACP.approxOr, ACP.approxOr_val}
\difficulty{1}
Fix a prime $p$, and let $P_1,\dots,P_{\mathrm{width}}$ be polynomials in the variables
$y_1,\dots,y_{\mathrm{vars}}$ over $\bbz/p$, together with subsets $S_1,\dots,S_\ell$ of
the index set $\{1,\dots,\mathrm{width}\}$. Then for every point $y :
\{1,\dots,\mathrm{vars}\} \to \bbz/p$, evaluating the randomized OR-approximating
polynomial $1 - \prod_{k=1}^{\ell}\bigl(1 - (\sum_{i \in S_k} P_i)^{p-1}\bigr)$ at $y$
yields the value obtained by feeding the scalars $P_1(y),\dots,P_{\mathrm{width}}(y)$
into the value-level OR approximator with the same subsets:
\[
\Bigl(1 - \prod_{k=1}^{\ell}\bigl(1 - (\textstyle\sum_{i \in S_k}
P_i)^{p-1}\bigr)\Bigr)\!\big|_{y}
\;=\;
1 - \prod_{k=1}^{\ell}\Bigl(1 - \bigl(\textstyle\sum_{i \in S_k}
P_i(y)\bigr)^{p-1}\Bigr).
\]
\end{lemma}

\begin{lemma}[Failure characterization of the value-level OR approximator]
\lean{ACP.approxOr_failure_iff}
\label{ACP.approxOr_failure_iff}
\leanok
\uses{ACP.OR_val, ACP.approxOr_val, ACP.one_sub_pow_card_sub_one}
\difficulty{4}
Fix a prime $p$, and let $v : \mathrm{Fin}\,\mathrm{width} \to \bbz/p$ be a vector
together with a family $S_1,\dots,S_\ell$ of subsets of $\{1,\dots,\mathrm{width}\}$.
Then the value-level OR approximator built from $v$ and $S$ disagrees with the
value-level OR detector of $v$ if and only if $v \neq 0$ yet $\sum_{i \in S_k} v_i = 0$
in $\bbz/p$ for every index $k$.
\end{lemma}

\begin{lemma}[Bad-seed count for a nonzero input]
\lean{ACP.count_bad_S}
\label{ACP.count_bad_S}
\leanok
\uses{ACP.OR_val, ACP.approxOr_failure_iff, ACP.approxOr_val}
\difficulty{5}
Fix a prime $p$ and let $v : \mathrm{Fin}\,\mathrm{width} \to \bbz/p$ be a nonzero
vector. Consider families $S = (S_1,\dots,S_\ell)$ of subsets $S_k \subseteq
\mathrm{Fin}\,\mathrm{width}$, indexed by $\mathrm{Fin}\,\ell$, and call such a family
*bad* for $v$ if the value-level OR approximator $1 - \prod_k\bigl(1 - (\sum_{i \in S_k}
v_i)^{p-1}\bigr)$ disagrees with the value-level OR detector $1 - \prod_k\bigl(1 -
v_k^{\,p-1}\bigr)$. Then the number of bad families, multiplied by $2^\ell$, is at most
the total number of families $S : \mathrm{Fin}\,\ell \to
\mathrm{Finset}(\mathrm{Fin}\,\mathrm{width})$.
\end{lemma}

\begin{lemma}[Cardinality of tuples of subsets]
\lean{ACP.approxSeed_card}
\label{ACP.approxSeed_card}
\leanok
\difficulty{2}
For natural numbers $\mathrm{width}$ and $\ell$, the number of $\ell$-tuples of subsets
of an $\mathrm{width}$-element set equals $2^{\mathrm{width}\cdot\ell}$; that is, the
set of functions assigning to each index $i \in \{0,\dots,\ell-1\}$ a subset of
$\{0,\dots,\mathrm{width}-1\}$ has exactly $2^{\mathrm{width}\cdot\ell}$ elements.
\end{lemma}

\begin{lemma}[Bad-seed bound for the value-level OR approximator]
\lean{ACP.count_bad_S_or}
\label{ACP.count_bad_S_or}
\leanok
\uses{ACP.OR_val, ACP.approxOr_failure_iff, ACP.approxOr_val, ACP.count_bad_S}
\difficulty{2}
Fix a prime $p$ and natural numbers $w$ and $\ell$, and let $v = (v_1,\dots,v_w)$ be an
arbitrary vector in $(\bbz/p)^w$ (no coordinate is assumed nonzero). Call a *seed* a
tuple $S = (S_1,\dots,S_\ell)$ of subsets of $\{1,\dots,w\}$, and call it *bad* when the
value-level OR approximator determined by $v$ and $S$ differs from the value-level OR
detector of $v$. Then, writing $N = (2^{w})^{\ell}$ for the total number of seeds, the
number of bad seeds satisfies
\[
(\text{number of bad seeds})\cdot 2^{\ell} \;\le\; N.
\]
\end{lemma}

\begin{definition}[List of OR-approximating polynomials]
\lean{ACP.approxOrPolyList}
\label{ACP.approxOrPolyList}
\leanok
\uses{ACP.approxOr}
The list of all polynomials $\texttt{ACP.approxOr}\,(P,S)$, one entry per random seed $S$. It
is a list rather than a set, so a polynomial arising from several seeds appears with the
corresponding multiplicity.
\end{definition}

\begin{lemma}[Number of OR-approximating polynomials]
\lean{ACP.approxOrPolyList_length}
\label{ACP.approxOrPolyList_length}
\leanok
\uses{ACP.approxOrPolyList, ACP.approxSeed_card}
\difficulty{2}
Fix a prime $p$ and, for a given width $\mathrm{width}$ and seed length $\ell$, a family
of polynomials $P_1,\dots,P_{\mathrm{width}}$ over $\bbz/p$ in a fixed number of
variables. Form the list whose entries are the OR-approximating polynomials, one for
each random seed — a seed being a tuple $(S_1,\dots,S_\ell)$ of $\ell$ subsets of the
index set $\{1,\dots,\mathrm{width}\}$, listed with multiplicity. Then this list has
length $2^{\mathrm{width}\cdot\ell}$, the number of such seeds.
\end{lemma}

\begin{theorem}[Few seeds miss the true OR at a point]
\lean{ACP.approxOr_pointwise_bad_count}
\label{ACP.approxOr_pointwise_bad_count}
\leanok
\uses{ACP.OR_val, ACP.approxOr, ACP.approxOr_eval_eq, ACP.approxSeed_card, ACP.count_bad_S_or}
\difficulty{2}
Fix a prime $p$, and let $P_1,\dots,P_{\mathrm{width}}$ be polynomials in
$\mathrm{vars}$ variables over $\bbz/p$, evaluated at a fixed point $y \in
(\bbz/p)^{\mathrm{vars}}$. A seed is a tuple $S = (S_1,\dots,S_\ell)$ of $\ell$ subsets
of $\{1,\dots,\mathrm{width}\}$, of which there are $2^{\mathrm{width}\cdot\ell}$ in
all; call $S$ bad at $y$ when the randomized OR-approximating polynomial built from the
$P_i$ and $S$, evaluated at $y$, differs from the value-level OR
\[
1 - \prod_{k=1}^{\mathrm{width}}\bigl(1 - P_k(y)^{\,p-1}\bigr).
\]
Then the number of seeds that are bad at $y$, multiplied by $2^\ell$, is at most
$2^{\mathrm{width}\cdot\ell}$; equivalently, at most a $2^{-\ell}$ fraction of all seeds
are bad at $y$.
\end{theorem}

\begin{theorem}[Existence of a good OR-approximating polynomial list]
\lean{ACP.exists_good_approxOr}
\label{ACP.exists_good_approxOr}
\leanok
\uses{ACP.approxOr, ACP.approxOrPolyList, ACP.approxOrPolyList_length, ACP.approxOr_pointwise_bad_count, ACP.approxOr_totalDegree}
\difficulty{3}
Fix a prime $p$ and natural numbers $\mathrm{vars}$, $\mathrm{width}$, and $\ell$, and
let $P_1,\dots,P_{\mathrm{width}}$ be polynomials in $\mathrm{vars}$ variables over
$\bbz/p$. Then there is a list $Ps$ of polynomials over $\bbz/p$ that is exactly the
list of OR-approximating polynomials — one entry for every seed $S$, a seed being an
assignment $k \mapsto S_k$ of a subset $S_k$ of the $\mathrm{width}$ positions to each
of the $\ell$ rounds — and this list has the following three properties. Its length is
$2^{\mathrm{width}\cdot\ell}$; every entry has total degree at most $(p-1)\,\ell\,\sup_i
\deg(P_i)$; and for every evaluation point $y \in (\bbz/p)^{\mathrm{vars}}$, the number
of seeds $S$ for which the OR-approximating polynomial of $S$ evaluated at $y$ differs
from the exact value $1 - \prod_{k=1}^{\mathrm{width}}\bigl(1 - P_k(y)^{p-1}\bigr)$,
when multiplied by $2^{\ell}$, is at most the length $2^{\mathrm{width}\cdot\ell}$ of
$Ps$.
\end{theorem}

\begin{definition}[Randomized AND-approximating polynomial]
\lean{ACP.approxAnd}
\label{ACP.approxAnd}
\leanok
\uses{ACP.approxOr}
The De Morgan dual of the OR approximator: $1 - \texttt{ACP.approxOr}\,(1 - P, S)$, where the
approximator is applied to the negated inputs $i \mapsto 1 - P_i$.
\end{definition}

\begin{theorem}[Total-degree bound for the AND approximator]
\lean{ACP.approxAnd_totalDegree}
\label{ACP.approxAnd_totalDegree}
\leanok
\uses{ACP.approxAnd, ACP.approxOr_totalDegree}
\difficulty{4}
Fix a prime $p$, and let $P_1,\dots,P_{\mathrm{width}}$ be polynomials over $\bbz/p$ in
the input variables, together with subsets $S_1,\dots,S_\ell$ of the index set
$\{1,\dots,\mathrm{width}\}$. Then the AND-approximating polynomial built from these
data has total degree at most
\[
(p-1)\,\ell\,\cdot\,\sup_{1 \le i \le \mathrm{width}} \deg(P_i).
\]
\end{theorem}

\begin{definition}[List of AND-approximating polynomials]
\lean{ACP.approxAndPolyList}
\label{ACP.approxAndPolyList}
\leanok
\uses{ACP.approxAnd}
The list of all polynomials $\texttt{ACP.approxAnd}\,(P,S)$, one entry per random seed $S$.
\end{definition}

\begin{lemma}[Length of the list of AND-approximating polynomials]
\lean{ACP.approxAndPolyList_length}
\label{ACP.approxAndPolyList_length}
\leanok
\uses{ACP.approxAndPolyList, ACP.approxSeed_card}
\difficulty{2}
Fix a prime $p$ and, for natural numbers $\mathrm{vars}$, $\mathrm{width}$, and $\ell$,
let $P_1,\dots,P_{\mathrm{width}}$ be polynomials in $\mathrm{vars}$ variables over
$\bbz/p$. A random seed is a tuple $(S_1,\dots,S_\ell)$ of subsets of the width
positions $\{1,\dots,\mathrm{width}\}$, of which there are $(2^{\mathrm{width}})^{\ell}
= 2^{\mathrm{width}\cdot\ell}$. The list obtained by recording the AND-approximating
polynomial of $P_1,\dots,P_{\mathrm{width}}$ for every such seed therefore has length
$2^{\mathrm{width}\cdot\ell}$.
\end{lemma}

\begin{theorem}[Pointwise bad-seed count for the AND approximator]
\lean{ACP.approxAnd_pointwise_bad_count}
\label{ACP.approxAnd_pointwise_bad_count}
\leanok
\uses{ACP.approxAnd, ACP.approxOr, ACP.approxOr_pointwise_bad_count}
\difficulty{4}
Fix a prime $p$, polynomials $P_1,\dots,P_{\mathrm{width}}$ over $\bbz/p$ in
$\mathrm{vars}$ variables, and an evaluation point $y \in (\bbz/p)^{\mathrm{vars}}$.
Consider the $2^{\mathrm{width}\cdot\ell}$ possible seeds $S = (S_1,\dots,S_\ell)$, each
$S_k$ a subset of $\{1,\dots,\mathrm{width}\}$, and let $N$ be the number of seeds for
which the randomized AND-approximating polynomial evaluated at $y$ differs from
$\prod_{k=1}^{\mathrm{width}}\bigl(1-(1-P_k(y))^{p-1}\bigr)$. Then
\[
N \cdot 2^{\ell} \le 2^{\mathrm{width}\cdot\ell}.
\]
\end{theorem}

\begin{theorem}[Existence of a good AND-approximating polynomial family]
\lean{ACP.exists_good_approxAnd}
\label{ACP.exists_good_approxAnd}
\leanok
\uses{ACP.approxAnd, ACP.approxAndPolyList, ACP.approxAndPolyList_length, ACP.approxAnd_pointwise_bad_count, ACP.approxAnd_totalDegree}
\difficulty{3}
Let $p$ be prime, let $\mathrm{vars}$, $\mathrm{width}$, and $\ell$ be natural numbers,
and let $P_1,\dots,P_{\mathrm{width}}$ be polynomials over $\bbz/p$ in $\mathrm{vars}$
variables. Call a *seed* an $\ell$-tuple $S=(S_1,\dots,S_\ell)$ of subsets of
$\{1,\dots,\mathrm{width}\}$, and for each seed $S$ write $A_S$ for the associated
AND-approximating polynomial; there are exactly $2^{\mathrm{width}\cdot\ell}$ seeds.
Then the list whose entries are the polynomials $A_S$, one per seed, has length
$2^{\mathrm{width}\cdot\ell}$, and every entry has total degree at most
$(p-1)\,\ell\,\sup_i \deg(P_i)$. Moreover, for every point
$y\in(\bbz/p)^{\mathrm{vars}}$, if $N(y)$ denotes the number of seeds $S$ for which
\[
A_S(y)\ \neq\ \prod_{k=1}^{\mathrm{width}}\Bigl(1-\bigl(1-P_k(y)\bigr)^{p-1}\Bigr)
\]
(the right-hand side being the exact conjunction value at $y$), then $N(y)\cdot
2^{\ell}\le 2^{\mathrm{width}\cdot\ell}$.
\end{theorem}

\begin{lemma}[{Case analysis on the $\mathrm{AC}^0[p]$ gate set}]
\lean{ACP.ACp_GateOps_cases}
\label{ACP.ACp_GateOps_cases}
\leanok
\uses{ACP.AC_GateOps, ACP.ACp_GateOps, ACP.FeedForward.GateOp.id, ACP.modGateOp}
\difficulty{3}
Let $p$ be a prime, and let $\mathrm{op}$ be a gate operation belonging to the
$\mathrm{AC}^0[p]$ gate set over the alphabet $\mathrm{Fin}\,2$. Then $\mathrm{op}$ is
one of exactly four kinds of gate: the identity gate $x \mapsto x_0$ on a single input,
the NOT gate $x \mapsto 1 - x_0$ on a single input, the unbounded fan-in AND gate $x
\mapsto \prod_{i} x_i$ for some arity $n$, or the unbounded $\mathrm{MOD}_p$ gate for
some arity $n$.
\end{lemma}

\begin{lemma}[Exact $\mathrm{MOD}_p$ on Boolean inputs]
\lean{ACP.exactMod_on_bits}
\label{ACP.exactMod_on_bits}
\leanok
\uses{ACP.bitify, ACP.cast_bitify_eq, ACP.modGateOp, ACP.one_sub_pow_card_sub_one}
\difficulty{3}
Let $p$ be a prime and let $\mathrm{inputs}_0, \dots, \mathrm{inputs}_{w-1}$ be elements
of $\mathbb{Z}/p$, each lying in $\{0,1\}$. Then
\[
1 - \Bigl(\sum_{i} \mathrm{inputs}_i\Bigr)^{p-1} =
\operatorname{MOD}_p\bigl(\mathrm{inputs}_0, \dots, \mathrm{inputs}_{w-1}\bigr),
\]
where the right-hand side is the output of the unbounded $\mathrm{MOD}_p$ gate on the
booleanizations of the inputs — equal to $1$ when $\sum_i \mathrm{inputs}_i = 0$ in
$\mathbb{Z}/p$ and to $0$ otherwise — regarded as an element of $\{0,1\} \subseteq
\mathbb{Z}/p$.
\end{lemma}

\begin{lemma}[Fermat power product computes Boolean AND]
\lean{ACP.exactAnd_on_bits}
\label{ACP.exactAnd_on_bits}
\leanok
\uses{ACP.bitify}
\difficulty{4}
Let $p$ be a prime, and let $x_0,\dots,x_{w-1}$ be elements of $\mathbb{Z}/p$, each
lying in $\{0,1\}$. Then
\[
\prod_i \bigl(1 - (1 - x_i)^{p-1}\bigr) = \prod_i \beta(x_i),
\]
where $\beta$ is the booleanization $\mathbb{Z}/p \to \mathrm{Fin}\,2$ (sending $1$ to
$1$ and everything else to $0$). The left-hand product is taken in $\mathbb{Z}/p$; on
the right, the product of the $\beta(x_i)$ is formed in $\mathrm{Fin}\,2$, and its
value—the natural number $0$ or $1$—is cast into $\mathbb{Z}/p$. Thus the Fermat-power
expression on the left evaluates to the Boolean AND of the inputs.
\end{lemma}

\begin{lemma}[{Low-degree approximating polynomials for an $\mathrm{AC}^0[p]$ gate}]
\lean{ACP.exists_poly_for_gate}
\label{ACP.exists_poly_for_gate}
\leanok
\uses{ACP.ACp_GateOps, ACP.ACp_GateOps_cases, ACP.approxAnd, ACP.approxAnd_pointwise_bad_count, ACP.approxAnd_totalDegree, ACP.approxSeed_card, ACP.bitify, ACP.cast_bitify_eq, ACP.exactAnd_on_bits, ACP.exactMod, ACP.exactMod_on_bits, ACP.exactMod_totalDegree, ACP.modGateOp}
\difficulty{6}
Let $p$ be a prime, and fix natural numbers $n$ and $\ell$. Let $\mathrm{op}$ be a gate
operation from the $\mathrm{AC}^0[p]$ gate set, and for each input wire $i$ of
$\mathrm{op}$ let $P_i$ be a polynomial in $\mathbb{Z}/p[X_1,\dots,X_n]$. Then there
exist a nonempty finite index set $\mathrm{Seed}$ with decidable equality and a family
of polynomials $(P_s)_{s\in\mathrm{Seed}}$ in $\mathbb{Z}/p[X_1,\dots,X_n]$ with the
following two properties. First, each $P_s$ has total degree at most
$(p-1)\,\ell\,\sup_i \deg(P_i)$. Second, for every Boolean assignment $x\in\{0,1\}^n$,
write $y\in(\mathbb{Z}/p)^n$ for the coordinatewise image of $x$ under the inclusion
$\mathrm{Fin}\,2\hookrightarrow\mathbb{Z}/p$; then, provided every value $P_i(y)$ lies
in $\{0,1\}$, the number of seeds $s\in\mathrm{Seed}$ for which $P_s(y)$ differs from
the gate output $\mathrm{op}\bigl(i\mapsto\mathrm{bitify}(P_i(y))\bigr)$, multiplied by
$2^\ell$, is at most $\abs{\mathrm{Seed}}$.
\end{lemma}
